\chapter{Register Rename Stage}

\section{Overview}

After an instruction has been decoded, it enters the Register Rename (IR) stage. The primary
objective of this stage is to eliminate false dependencies caused by reuse of architectural
register names.

In the Sargantana processor, the rename stage consists of three major components:
\begin{itemize}
  \item Register Alias Table (Rename Table)
  \item Free List
  \item Instruction Queue
\end{itemize}

These modules are implemented in:

\begin{center}
\rtl{rtl/datapath/rtl/ir\_stage/}
\end{center}

Important RTL files include:

\begin{rtlrefs}
  \rtlitem{rtl/datapath/rtl/ir_stage/rtl/rename_table.sv}
  \rtlitem{rtl/datapath/rtl/ir_stage/rtl/free_list.sv}
  \rtlitem{rtl/datapath/rtl/ir_stage/rtl/instruction_queue.sv}
\end{rtlrefs}

\begin{figure}[H]
  \centering
  \includegraphics[width=0.33\textwidth]{ch07_rename_overview.png}
  \caption{Rename stage: Instruction Queue, Renaming Tables, Free Lists, and Recovery Logic}
  \label{fig:rename-overview}
\end{figure}

\section{Register Renaming}

Modern superscalar processors allow multiple instructions to execute simultaneously. If two
instructions write to the same architectural register, false dependencies arise even when the
instructions are logically independent.

Register renaming solves this problem by mapping architectural registers onto a much larger
pool of physical registers.

For example:
\begin{verbatim}
ADD x5, x1, x2
SUB x5, x3, x4
\end{verbatim}

Architecturally, both instructions write to \texttt{x5}.

\textbf{Without renaming:}
\begin{verbatim}
ADD -----> x5
SUB -----> x5
\end{verbatim}
The second instruction must wait.

\textbf{With renaming:}
\begin{verbatim}
ADD -----> P17
SUB -----> P22
\end{verbatim}
Both instructions can execute independently.

This mapping is maintained inside:

\begin{center}
\rtl{rtl/datapath/rtl/ir\_stage/rtl/rename\_table.sv}
\end{center}

\section{Register Alias Table (RAT)}

The Register Alias Table (RAT), implemented as the Rename Table, stores the current mapping
between architectural registers and physical registers.

Rather than reading an architectural register directly, later stages first consult the Rename Table
to determine which physical register currently holds the latest value.

Whenever a new destination register is renamed, the table updates its mapping.

The implementation resides in:

\begin{center}
\rtl{rtl/datapath/rtl/ir\_stage/rtl/rename\_table.sv}
\end{center}

The Rename Table therefore becomes the central structure responsible for maintaining correct
register mappings throughout speculative execution.

\section{Physical Register Allocation}

Each new destination register requires a free physical register.

Instead of searching all registers every cycle, Sargantana maintains a dedicated Free List.

The Free List keeps track of unused physical registers.

Whenever an instruction requires a destination register, the allocation proceeds as shown below.

\begin{figure}[H]
  \centering
  \includegraphics[width=0.3\textwidth]{ch07_physical_register_allocation.png}
  \caption{Physical register allocation flow: Instruction arrives $\rightarrow$ Request Physical Register $\rightarrow$ Free List selects register $\rightarrow$ Rename Table updated}
  \label{fig:phys-reg-alloc}
\end{figure}

When an instruction finally commits, its old physical register is returned to the Free List for
reuse.

This functionality is implemented in:

\begin{center}
\rtl{rtl/datapath/rtl/ir\_stage/rtl/free\_list.sv}
\end{center}

Together, the Free List and Rename Table enable efficient dynamic register allocation.

\section{Hazard Elimination}

The rename stage primarily eliminates false dependencies:
\begin{itemize}
  \item Write After Write (WAW)
  \item Write After Read (WAR)
\end{itemize}

Because each destination register receives a unique physical register, later instructions no
longer overwrite the results of earlier instructions.

Only true data dependencies (RAW) remain.

The renamed instructions are temporarily stored inside the Instruction Queue before being
issued to the Register Read stage.

Instruction Queue implementation:

\begin{center}
\rtl{rtl/datapath/rtl/ir\_stage/rtl/instruction\_queue.sv}
\end{center}

The rename stage therefore prepares instructions for high-performance out-of-order execution
while maintaining correct architectural state.
